\subsection{One and two terms}\label{subsec:baselines-lift}

With two odd powers, an energy comparison orders any two putative cycles.
For attainment it suffices to perturb a single family of circles.

\begin{proof}[Proof of Theorem~\ref{thm:lower}\textup{(i)}]
An even $F$ has identical half-profiles. Every periodic orbit belongs to
an open annulus of periodic orbits by transversality of its two axis hits. This includes two even powers.
For a single odd monomial the divergence $-F'(x)$ has one strict sign off
$x=0$, so Green's theorem excludes a periodic orbit. This proves the
one-term assertion. With one odd and
one even power, $F(r)-F(-r)$ has a strict sign on $r>0$. A periodic
orbit would supply two transverse returns from its upper intercept.
Article, Lemma~2.1(ii) orders their lower hits strictly.

It remains to bound the number for two odd powers. Equal coefficient signs
again give divergence of one sign. With opposite signs we may reverse
time and replace $y$ by $-y$ so that
\[
 F(x)=-c x^{2p+1}+e x^{2q+1},\qquad c,e>0,\quad 0\le p<q .
\]
If $p=0$, the damping $F'$ is negative at zero and increases on
the positive half-axis, so Massera's uniqueness theorem applies
\cite[pp.~367--369]{Massera}. For $p>0$, $F'$ initially decreases.
The energy comparison below covers both cases. Let $a=(c/e)^{1/(2q-2p)}$. Then $F<0$ on $(0,a)$ and $F$ increases
strictly from zero on $[a,\infty)$. Central reflection preserves the
field's orbits. Each closed orbit is centrally symmetric, since its image
is either itself or a strictly nested closed orbit, and applying the
reflection twice rules out strict nesting.

Write $V=(x^2+y^2)/2$. Along an orbit $\dot V=-xF(x)$ and
$\dd V=F(x)\dd y$. On the right half of a closed orbit $V$ takes the
same value at the two axis hits. This arc must pass $x=a$, since $V$
strictly increases while $0<x<a$. Denote the upper and lower ordinates
at abscissa $x$ by $y_o(x)>F(x)>y_i(x)$. The gain of $V$ in $0<x<a$ is
\[
 \int_0^a x[-F(x)]\!
 \left(\frac1{y_o(x)-F(x)}+\frac1{F(x)-y_i(x)}\right)\dd x .
\]
For a strictly enclosing orbit the gain is smaller, since both vertical
gaps in the denominators increase. The loss beyond $a$ is
$\int_{y_i(a)}^{y_o(a)} F(x(y))\dd y$, where the right-hand arc is a
single graph over $y$. On the enclosing arc the interval of integration
is larger and $x(y)$ is strictly larger on the common interval.
As $F$ is strictly increasing there, the loss is larger. Two nested
cycles cannot both balance their gain and loss. This proves
$\LieM(2)\le1$, including $p>0$.

For attainment with a hyperbolic cycle take $F(x)=\delta(x^3-x)$.
On a compact annulus about radius $R_*=2/\sqrt3$, the return after one
clockwise revolution is smooth in $R$ and $\delta$. Along the unperturbed
circle $x=R\cos t$, integration gives
\[
 V_{\rm return}-V_{\rm initial}
  =\delta\!\int_0^{2\pi}(R^2\cos^2t-R^4\cos^4t)\dd t+O(\delta^2)
  =\pi\delta R^2(1-3R^2/4)+O(\delta^2)
\]
in $C^1$ on that annulus. Divide by $\delta$ and apply the implicit
function theorem at the simple positive zero $R_*$. For all sufficiently
small $\delta>0$ the resulting orbit is hyperbolic, proving $H(2)=1$.
\end{proof}
